\section{What \sgt{} Commits To}
\label{sec:design}

\begin{figure*}[t]
  \centering
  % Figure 3. The whole model in one picture: what a save writes, what a version
% is, and where the bytes in your files come from. Same session as Figure 1.
\newcommand{\FigThreeCaption}{What one save writes down, what a version is, and
where the bytes in a file come from. A save writes
one record per changed function and files the records under a guess at the request
they served, which the developer corrects with one command (left). A version is
the set of those records, so taking out the price cache is subtracting six records
from eighteen, and two rules decide which sets are allowed (middle). The files are
rebuilt from the set with the text between functions kept exactly, so the same set
always produces the same bytes (right). Two kinds of mark in the right-hand column
are the two consequences \sgt{} reports about this revert and does not act on. The hollow
triangle on \protect\fsym{api.py::handle} means the retry policy edited that
function after the cache did, so \sgt{} took the cache's lines back out of the
current text and the developer's test suite decides whether the result stands. The
open circles are three functions that still call
\protect\fsym{cache.py::get\_cached}, which \sgt{} names and leaves alone.
\protect\fsym{db.py::query} is the one function that goes back to text it had
earlier in the day.}

\newcommand{\FigThreeDescription}{Three panels. The left panel shows a terminal
transcript of one save command, followed by the three feature groups the save
landed in, one of them newly created and unnamed, and then a card listing the
five fields one edit record holds. The middle panel lists the fourteen functions
this session changed, grouped by file, with the edit marks for each on a short
rail, in two columns. The left column is the set as it stands and the right
column is the same list after the price cache has been reverted, where the cache
edits are drawn as dotted outlines, one file is marked as removed, one function is
marked with a hollow triangle, and three functions are marked with an open circle.
The right panel draws the file
api.py twice as a vertical stack of bands, once before and once after, where
coloured bands are function bodies and grey bands are the text between them,
which is identical in both.}

\begin{tikzpicture}[x=1mm, y=1mm, line cap=round, font=\sffamily]

\node[hd, anchor=west] at (0,78)   {1\quad WHAT A SAVE WRITES};
\node[hd, anchor=west] at (52,78)  {2\quad WHAT A VERSION IS};
\node[hd, anchor=west] at (117,78) {3\quad WHAT IS IN YOUR FILES};
\draw[color=sgtGrey!40, line width=0.3pt] (0,76.5) -- (176,76.5);
\draw[color=sgtGrey!30, line width=0.3pt] (49.5,11) -- (49.5,76.5);
\draw[color=sgtGrey!30, line width=0.3pt] (114,11) -- (114,76.5);
\draw[color=sgtGrey!40, line width=0.3pt] (0,9.6) -- (176,9.6);

% ==================== panel 1: what a save writes =========================
\node[anchor=north west, align=left, inner sep=0, text=sgtInk,
      font=\ttfamily\fontsize{5.2}{7}\selectfont] at (0,74.6)
  {\$ sgt save -m "retry the upstream call\\
   \hphantom{\$ sgt save -m "}when it returns a 5xx"\\
   \ \ save 7f10de\ \ 6 edits};

\oprew{1.3}{62.2}{fRetry}
\node[anchor=north west, align=left, inner sep=0, text=sgtInk,
      font=\ttfamily\fontsize{5.2}{6.6}\selectfont] at (3.4,63.2)
  {\textcolor{fRetry}{new feature} (7c02b4d1)\ \ 3 edits\\
   retry.py::Retry\ \ retry.py::\_backoff\\
   api.py::handle\\
   unnamed. name it with\\
   sgt feature rename 7c02b4d1 "retry policy"};

\oprew{1.3}{45.6}{fCache}
\node[anchor=north west, align=left, inner sep=0, text=sgtInk,
      font=\ttfamily\fontsize{5.2}{6.6}\selectfont] at (3.4,46.6)
  {\textcolor{fCache}{price cache} (001f73c2)\ \ 2 edits\\
   cache.py::get\_cached\ \ cache.py::\_ttl};

\oprew{1.3}{39.4}{sgtGrey}
\node[anchor=north west, align=left, inner sep=0, text=sgtInk,
      font=\ttfamily\fontsize{5.2}{6.6}\selectfont] at (3.4,40.4)
  {\textcolor{sgtGrey}{row accessors} (9d20b18a)\ \ 1 edit\\
   util.py::audit};

% one record, in full
\draw[card, color=sgtGrey!45, fill=sgtPaper] (0,14.6) rectangle (47,34.2);
\fill[fRetry] (0,31.6) rectangle (47,34.2);
\node[anchor=west, inner sep=0, text=sgtPaper,
      font=\ttfamily\fontsize{5.2}{6}\selectfont] at (1.2,32.9) {edit 4b1c9f28};
\node[anchor=north west, align=left, inner sep=0, text=sgtInk,
      font=\ttfamily\fontsize{5.2}{6.9}\selectfont] at (1.2,30.6)
  {function\ \ \ \ api.py::handle\\
   kind\ \ \ \ \ \ \ \ extended\\
   started at\ \ 9ac7e1\\
   produced\ \ \ \ b41c08\\
   built on\ \ \ \ retry.py::Retry};

% ==================== panel 2: what a version is ==========================
\node[anchor=north west, align=left, text width=61mm, inner sep=0, text=sgtInk!85,
      font=\sffamily\fontsize{5.6}{6.6}\selectfont] at (52,74.8)
  {\textbf{rule 1}\ \ no edit without the edits it was built on.\ \ \ \
   \textbf{rule 2}\ \ one live version of each function.};

\node[anchor=north west, align=left, inner sep=0, text=sgtInk,
      font=\sffamily\bfseries\fontsize{5.6}{6.4}\selectfont] at (52,68.6)
  {THE SET NOW};
\node[anchor=north west, inner sep=0, text=sgtGrey,
      font=\sffamily\fontsize{5.2}{6}\selectfont] at (52,66.0)
  {18 edits, 14 functions};
\node[anchor=north west, align=left, inner sep=0, text=sgtInk,
      font=\ttfamily\fontsize{5.4}{6.2}\selectfont] at (84,68.6)
  {\$ sgt revert "price cache"};
\node[anchor=north west, inner sep=0, text=sgtGrey,
      font=\sffamily\fontsize{5.2}{6}\selectfont] at (84,66.0)
  {12 edits, 12 functions};

% #1 x, #2 y, #3 file name
\newcommand{\FTfile}[3]{\node[anchor=north west, inner sep=0, text=sgtInk!65,
  font=\sffamily\bfseries\fontsize{5.2}{6}\selectfont] at (#1,#2) {#3};}
% #1 x, #2 y, #3 symbol name, #4 colour
\newcommand{\FTsym}[4]{\node[anchor=north west, inner sep=0, text=#4,
  font=\ttfamily\fontsize{5.2}{6}\selectfont] at (#1,#2) {#3};}
% a hollow up-triangle: this version was re-drafted, and the tests decide
\newcommand{\FTredraft}[2]{\draw[sgtInk!85, line width=0.4pt, fill=sgtPaper]
  ({#1-0.95},{#2-0.85}) -- ({#1+0.95},{#2-0.85}) -- (#1,{#2+1.0}) -- cycle;}
% an open circle: this function still calls something the revert removed
\newcommand{\FTcalls}[2]{\draw[sgtInk!85, line width=0.4pt, fill=sgtPaper]
  (#1,#2) circle (0.85);}
% #1 x of rail start, #2 y of rail
\newcommand{\FTrail}[3]{\draw[rail=#3] (#1,#2) -- ({#1+9},#2);}
\newcommand{\FTraildim}[2]{\draw[raildim] (#1,#2) -- ({#1+9},#2);}

% ---- left column: the set as it stands
\FTfile{52}{63.6}{api.py}
\FTsym{54}{60.8}{handle}{sgtInk}
  \FTrail{69.5}{59.9}{fRetry}
  \opext{70.6}{59.9}{fRate}\opext{74.0}{59.9}{fCache}\opext{77.4}{59.9}{fRetry}
\FTsym{54}{58.0}{\_bucket}{sgtInk}
  \FTrail{69.5}{57.1}{fRate}\opadd{70.6}{57.1}{fRate}
\FTsym{54}{55.2}{\_key}{sgtInk}
  \FTrail{69.5}{54.3}{fRate}\opadd{70.6}{54.3}{fRate}
\FTsym{54}{52.4}{\_window}{sgtInk}
  \FTrail{69.5}{51.5}{fRate}\opadd{70.6}{51.5}{fRate}
\FTsym{54}{49.6}{export\_row}{sgtInk}
  \FTrail{69.5}{48.7}{sgtGrey}\opext{70.6}{48.7}{sgtGrey}
\FTfile{52}{46.8}{cache.py}
\FTsym{54}{44.0}{get\_cached}{sgtInk}
  \FTrail{69.5}{43.1}{fCache}
  \opadd{70.6}{43.1}{fCache}\oprew{74.0}{43.1}{fCache}
\FTsym{54}{41.2}{\_ttl}{sgtInk}
  \FTrail{69.5}{40.3}{fCache}
  \opadd{70.6}{40.3}{fCache}\opext{74.0}{40.3}{fCache}
\FTfile{52}{38.4}{db.py}
\FTsym{54}{35.6}{query}{sgtInk}
  \FTrail{69.5}{34.7}{fCache}\opext{70.6}{34.7}{fCache}
\FTsym{54}{32.8}{fetch\_row}{sgtInk}
  \FTrail{69.5}{31.9}{sgtGrey}\opmove{70.6}{31.9}{sgtGrey}
\FTfile{52}{30.0}{retry.py}
\FTsym{54}{27.2}{Retry}{sgtInk}
  \FTrail{69.5}{26.3}{fRetry}\opadd{70.6}{26.3}{fRetry}
\FTsym{54}{24.4}{\_backoff}{sgtInk}
  \FTrail{69.5}{23.5}{fRetry}\opadd{70.6}{23.5}{fRetry}
\FTfile{52}{21.6}{util.py}
\FTsym{54}{18.8}{load\_row}{sgtInk}
  \FTrail{69.5}{17.9}{sgtGrey}\opext{70.6}{17.9}{sgtGrey}
\FTsym{54}{16.0}{save\_row}{sgtInk}
  \FTrail{69.5}{15.1}{sgtGrey}\opext{70.6}{15.1}{sgtGrey}
\FTsym{54}{13.2}{audit}{sgtInk}
  \FTrail{69.5}{12.3}{sgtGrey}\opext{70.6}{12.3}{sgtGrey}

% ---- right column: the same list, six edits subtracted
\FTfile{84}{63.6}{api.py}
\FTsym{86}{60.8}{handle}{sgtInk}
  \FTrail{101.5}{59.9}{fRetry}
  \opext{102.6}{59.9}{fRate}\opghost{106.0}{59.9}{fCache}\opext{109.4}{59.9}{fRetry}
  \FTredraft{112.6}{59.9}
\FTsym{86}{58.0}{\_bucket}{sgtInk}
  \FTrail{101.5}{57.1}{fRate}\opadd{102.6}{57.1}{fRate}
\FTsym{86}{55.2}{\_key}{sgtInk}
  \FTrail{101.5}{54.3}{fRate}\opadd{102.6}{54.3}{fRate}
\FTsym{86}{52.4}{\_window}{sgtInk}
  \FTrail{101.5}{51.5}{fRate}\opadd{102.6}{51.5}{fRate}
\FTsym{86}{49.6}{export\_row}{sgtInk}
  \FTrail{101.5}{48.7}{sgtGrey}\opext{102.6}{48.7}{sgtGrey}
  \FTcalls{112.6}{48.7}
\FTfile{84}{46.8}{cache.py}
\node[anchor=east, inner sep=0, text=sgtGrey,
      font=\sffamily\itshape\fontsize{5}{5.8}\selectfont] at (113.5,46.0)
  {file removed};
\FTsym{86}{44.0}{get\_cached}{sgtGrey!80}
  \FTraildim{101.5}{43.1}
  \opghost{102.6}{43.1}{fCache}\opghost{106.0}{43.1}{fCache}
\FTsym{86}{41.2}{\_ttl}{sgtGrey!80}
  \FTraildim{101.5}{40.3}
  \opghost{102.6}{40.3}{fCache}\opghost{106.0}{40.3}{fCache}
\FTfile{84}{38.4}{db.py}
\FTsym{86}{35.6}{query}{sgtInk}
  \draw[raildim] (101.5,34.7) -- (104.0,34.7);\opghost{102.6}{34.7}{fCache}
  \node[anchor=east, inner sep=0, text=sgtGrey,
        font=\sffamily\itshape\fontsize{5}{5.8}\selectfont] at (113.5,34.7)
    {back to 14:02};
\FTsym{86}{32.8}{fetch\_row}{sgtInk}
  \FTrail{101.5}{31.9}{sgtGrey}\opmove{102.6}{31.9}{sgtGrey}
\FTfile{84}{30.0}{retry.py}
\FTsym{86}{27.2}{Retry}{sgtInk}
  \FTrail{101.5}{26.3}{fRetry}\opadd{102.6}{26.3}{fRetry}
\FTsym{86}{24.4}{\_backoff}{sgtInk}
  \FTrail{101.5}{23.5}{fRetry}\opadd{102.6}{23.5}{fRetry}
\FTfile{84}{21.6}{util.py}
\FTsym{86}{18.8}{load\_row}{sgtInk}
  \FTrail{101.5}{17.9}{sgtGrey}\opext{102.6}{17.9}{sgtGrey}
  \FTcalls{112.6}{17.9}
\FTsym{86}{16.0}{save\_row}{sgtInk}
  \FTrail{101.5}{15.1}{sgtGrey}\opext{102.6}{15.1}{sgtGrey}
  \FTcalls{112.6}{15.1}
\FTsym{86}{13.2}{audit}{sgtInk}
  \FTrail{101.5}{12.3}{sgtGrey}\opext{102.6}{12.3}{sgtGrey}

% ==================== panel 3: what is in your files ======================
\node[anchor=north west, text width=59mm, align=left, inner sep=0, text=sgtInk!85,
      font=\sffamily\fontsize{5.6}{6.6}\selectfont] at (117,74.8)
  {\sgt{} writes a file by putting each live version back in place and keeping
   every character between them.};

% #1 x left, #2 y bottom, #3 height, #4 colour, #5 label
\newcommand{\FTband}[5]{%
  \fill[#4!30, rounded corners=0.4mm] ({#1+0.6},#2) rectangle ({#1+25.4},{#2+#3});
  \node[anchor=west, inner sep=0, text=sgtInk!85,
        font=\ttfamily\fontsize{5}{5.8}\selectfont] at ({#1+2.2},{#2+#3/2}) {#5};}
% #1 x left, #2 y bottom, #3 height: text sgt keeps verbatim
\newcommand{\FTgap}[3]{%
  \fill[sgtGrey!12] ({#1+0.6},#2) rectangle ({#1+25.4},{#2+#3});
  \foreach \i in {0,1,2} {\fill[sgtGrey!45] ({#1+2.2},{#2+#3-1.1-\i*1.1})
     rectangle ({#1+2.2+13-\i*3.2},{#2+#3-0.65-\i*1.1});}}

\node[anchor=south west, inner sep=0, text=sgtInk,
      font=\ttfamily\fontsize{5.4}{6}\selectfont] at (117,68.6) {api.py};
\node[anchor=south west, inner sep=0, text=sgtInk,
      font=\ttfamily\fontsize{5.4}{6}\selectfont] at (148,68.6) {api.py};
\node[anchor=south east, inner sep=0, text=sgtGrey,
      font=\sffamily\fontsize{5}{5.8}\selectfont] at (143,68.6) {before};
\node[anchor=south east, inner sep=0, text=sgtGrey,
      font=\sffamily\fontsize{5}{5.8}\selectfont] at (174,68.6) {after};

\foreach \x in {117,148} {
  \draw[card, color=sgtGrey!45, fill=sgtPaper] (\x,20.5) rectangle ({\x+26},67.5);
  \FTgap{\x}{63.6}{3.4}
  \FTband{\x}{58.4}{4.6}{fRate}{\_bucket}
  \FTband{\x}{53.4}{4.6}{fRate}{\_key}
  \FTband{\x}{48.4}{4.6}{fRate}{\_window}
  \FTgap{\x}{44.6}{3.4}
  \FTgap{\x}{30.0}{3.4}
  \FTband{\x}{21.4}{8.2}{sgtGrey}{export\_row}
}
% handle: one version, with a bar saying which requests contributed to it
\FTband{117}{34.0}{10.2}{fRetry}{handle}
\fill[fRate]  (117.6,34.0) rectangle (118.7,37.4);
\fill[fCache] (117.6,37.4) rectangle (118.7,40.8);
\fill[fRetry] (117.6,40.8) rectangle (118.7,44.2);
\FTband{148}{34.0}{10.2}{fRetry}{handle}
\fill[fRate]  (148.6,34.0) rectangle (149.7,39.1);
\fill[fRetry] (148.6,39.1) rectangle (149.7,44.2);
\FTredraft{171.6}{39.1}

\node[anchor=north west, text width=59mm, align=left, inner sep=0, text=sgtInk!85,
      font=\sffamily\fontsize{5.6}{6.6}\selectfont] at (117,19.0)
  {The grey bands are the imports, the comments, and the blank lines, which \sgt{}
   keeps and never rewrites. The bar left of \fsym{handle} names the requests that
   contributed to the version in the file.};

% ==================== bottom: what each one buys ==========================
\node[anchor=north west, text width=47mm, align=left, inner sep=0, text=sgtInk,
      font=\sffamily\fontsize{5.8}{6.8}\selectfont] at (0,8.4)
  {Every field was read out of the files, and the sentence at the top is all the
   developer typed. No diff was sorted by hand at any point.};
\node[anchor=north west, text width=61mm, align=left, inner sep=0, text=sgtInk,
      font=\sffamily\fontsize{5.8}{6.8}\selectfont] at (52,8.4)
  {Each of \cmd{revert}, \cmd{restore}, \cmd{diff} and \cmd{switch} is one
   subtraction or one union over this list, so \sgt{} can show the resulting files
   before writing any of them.};
\node[anchor=north west, text width=59mm, align=left, inner sep=0, text=sgtInk,
      font=\sffamily\fontsize{5.8}{6.8}\selectfont] at (117,8.4)
  {A file is a rebuild from a set. Revert the cache twice and restore it once and
   the bytes match the file that was there before any of it.};

\end{tikzpicture}

  \Description{\FigThreeDescription}
  \caption{\FigThreeCaption}
  \label{fig:model}
\end{figure*}

A version control system's unit of record decides which questions a developer can
ask of their own history. The unit should match the grain a developer thinks
at---because those are the questions they will actually ask. \sgt{} records one edit per changed function
and files it under the request the edit was made for. This section states that
commitment and the three that follow from it, each with the price it carries.
Figure~\ref{fig:model} shows the first two of them at work on the sitting from
Figure~\ref{fig:two-records}: what one save writes down, and what a version is.

The design philosophy is that a tool should refuse rather than guess wrong, and
should report rather than act silently. A differently-minded designer would merge
concurrent edits to one function the way git merges concurrent edits to one file,
on the grounds that most such merges produce working code. We tried that first and
abandoned it because when the merge fails the developer has lost work they did not
know was at risk, and the failure appears somewhere other than where the merge
happened. Another reasonable design would ask the developer at save time which
request each function belongs to. We tried that too, and the developer cannot
answer it without reading the diff, which is the labour the tool was supposed to
eliminate. The four commitments below are what survived: each one is the design we
arrived at after building the alternative and finding its cost.

\subsection{One edit to one function, filed under the request}
\label{sec:grain}

A save reads the working files, compares them against the last recorded state, and
writes one record per changed function, naming the function, the version of it the
edit started from, the version it produced, and those other functions the new code
refers to that \sgt{} can identify without guessing: a reference is recorded when the
name it uses resolves to exactly one function in the codebase, and left unrecorded
when the same short name is defined in several places, because a reference attributed
to the wrong definition is worse than a missing one. The developer supplies one sentence,
\cmd{sgt save -m "retry the upstream call when it returns a 5xx"}, and \sgt{} reads
everything else out of the files, so a record can be no more wrong than the code it
was read from.

We chose the function as the unit because the three alternatives each fail on the
same ordinary case. The price cache in Figure~\ref{fig:model} added a parameter to
\sym{db.py::query} and passed a matching argument from \sym{api.py::handle}, and
those two changes share no line and no file, so a record that bottoms out at a line
has nowhere to write down that they were made for one reason, and the developer who
removes the cache a week later learns about the second one when the import fails. A
record whose unit is the file cannot say that two changes to \sym{api.py} served
different purposes, which is the default case in Figure~\ref{fig:two-records} where
three requests all edit the same handler. GitButler uses this unit and solves the
problem it can solve, which is the case where one feature per file
suffices~\cite{gitbutler2026}. A record whose unit is the commit holds whatever was
pending when the developer typed the command, which in
Figure~\ref{fig:two-records} is three requests and a rename the agent did on its
own initiative, divided across two commits by the clock. A function is the smallest
thing that both the developer and the agent name out loud while the work is
happening, and the feature location literature works at that same unit because the
code serving one concern is a set of functions scattered across
files~\cite{biggerstaff1993,dit2013}.

\sgt{} finds functions in Python, TypeScript, and TSX with
tree-sitter~\cite{treesitter} and records every other file as its whole content, at
the grain git already offers, so the rest of this section applies only to files
\sgt{} can parse. Within the three parsed languages a function can be smaller than a purpose. The
rename in Figure~\ref{fig:two-records} moved \sym{db.py::fetch} to
\sym{fetch\_row} and updated four call sites, so \sgt{} writes five records for
what the developer thinks of as one action. Because those five records share one
name, the mismatch is invisible until somebody reads the list underneath it. A function can also be
larger than a purpose, which happens when two requests change opposite ends of one
fifty line function, and Section~\ref{sec:fork} describes what \sgt{} does then.

\subsection{A version is a set of edits, and the files are rebuilt from it}
\label{sec:sets}

Two rules decide which sets are versions. No edit is included without the edits it
was built on, and each function has one included version. \sgt{} rebuilds each file
by writing out the included version of every function it contains and keeping the
text between functions exactly as it was, which produces the bytes that are checked
out. That interstitial text (whitespace, comments between definitions) is itself
recorded as a layout symbol so that its position is deterministic under
set operations. (Section~\ref{sec:layout-seam} reports the cost of that choice.)
The middle panel of Figure~\ref{fig:model} is one such set before and after a
revert, and the right panel is the file each set produces.

The alternative is to store patches and replay them, which is what Darcs and Pijul
do at the line level~\cite{roundy2005,pijul2024}. We chose sets over replay
because replay introduces ordering: reverting a feature means computing an inverse
patch and applying it to the current state, and the result depends on what the
current state is. A set-based revert is idempotent (removing what is already absent
changes nothing) and its inverse is exact (restoring a removed set reproduces the
original bytes with no drift). Those two properties make the preview trustworthy,
because the files the preview shows are the files the operation will produce,
regardless of what else has been recorded since. The cost is that every function
must have exactly one included version at any moment, which is the constraint the
fork rule (Section~\ref{sec:fork}) enforces.

Because a version is a set, \cmd{revert}, \cmd{restore}, \cmd{diff}, and
\cmd{switch} all take a set and return a set, and two properties follow that a
developer can check on their own repository. Reverting the price cache and then
restoring it produces a file byte-identical to the one that was there before either
command ran, because both operations are set arithmetic over the same records and
neither reapplies a patch. Each of the four can also be shown in full before it
happens, since working out the resulting files takes no more than rebuilding them
from the proposed set: \cmd{sgt revert "price cache"} draws the twelve edits that
would remain, names the functions whose text would change, and asks
\cmd{apply this revert? [y/N]}. Run with no terminal attached it prints the same
preview and declines, so a revert nobody watched does not happen. The safety
property (no mutation without confirmation) held throughout the evaluation; the
\emph{information} property did not. During the evaluation the machine-readable
preview omitted the subtraction report entirely and its one-line summary asserted
``nothing depends on it'' while its own payload carried a non-zero dependent
count, so an agent that checked the preview to decide whether to confirm was given
less information than a developer reading the terminal.

Three requests wrote \sym{api.py::handle}, and in Figure~\ref{fig:model} it carries
an edit from the rate limiter, one from the cache, and one from the retry policy, in
that order. The retry policy was built on the cache's edit (it starts from the
version the cache left behind), so the dependency rule chains them: removing the
cache alone would orphan the retry. We shipped that strict interpretation first---a
revert of the cache also removed the retry---and it emptied a test repository,
because on any history where features interleave inside a shared function, removing
an early edit takes every later edit with it. The default now keeps the later edits
and subtracts the reverted feature's contribution from the current text of the
shared function with a three-way merge, the same operation git performs when it
merges a text file. The version of
\sym{api.py::handle} that Figure~\ref{fig:model} leaves in the file is produced that
way, and the hollow diamond beside it is how \sgt{} says so. The old behaviour
remains available as \cmd{revert --take-dependents}, and its help text says what it
removes.

That three-way merge is the only place \sgt{} changes text that no person wrote, so
\sgt{} names every function it changes that way and every function it decided to
leave alone. When the subtraction overlaps a line a later edit changed, \sgt{} leaves
the function exactly as it is and prints \cmd{kept unchanged (the removal overlaps
later edits\ttdash needs your edit)} with the function named, and when a function
\sgt{} is keeping still calls a function the revert removed, it prints
\cmd{still references removed code (fix or revert separately)}. Both are reports and neither stops the revert, because the alternative is a
command that refuses to finish over a consequence the developer may already
intend. The second warning depends on the reference field described in
Section~\ref{sec:grain}, which requires a globally unique function name, so on
repositories where short names recur (monorepos, projects with many modules) the
warning is structurally absent and the developer's test suite is the only oracle.
Across the evaluation corpus, 6.7\% of records carry any reference edge (the
per-repository range is 0.0--20.5\%), so the dependency-based refusal in
Section~\ref{sec:walkthrough}'s Wednesday scenario represents a capability most
records cannot trigger; a repository whose function names are not globally unique
sees none of it. The cost of the default is that after the merge the shared function is text the
developer has not read, and their test suite decides whether the revert was
correct. Naming the function before the revert runs is as far as we are willing
to go on the developer's behalf.

\subsection{Two versions of one function, and a person decides}
\label{sec:fork}

\begin{figure}[t]
  \centering
  % The fork rule: one function, two versions, and who decides. Prints as
% Figure 4, single column, beside Section 4.3. Same panels, glyphs and palette
% as Figures 1 to 3. The instance is the one the walkthrough reaches in 5.5.
\newcommand{\FigForkCaption}{Two people change \protect\fsym{api.py::handle}
from the same saved version, one cutting three lines from the top of it and one
adding a line forty lines below. Git merges the two and prints nothing, and
nobody has run the function in the state the merge produces. \sgt{} records both
versions, includes neither, and names the function. Of everything \sgt{} refuses,
only this refusal can stop work that would have been correct.}

\newcommand{\FigForkDescription}{Three stacked panels. The top panel shows one
recorded version of a function splitting into two, one change removing three
lines at the top of the function and one adding a line at the bottom, each
marked with a small figure of a person, with the two resulting function bodies
drawn side by side. The middle panel shows git's result, a single function body
with the top lines gone and the bottom line added, annotated as a state nobody
ran. The bottom panel shows sgt's result, both versions drawn as dotted outlines
to mean neither is included, and one solid version after the developer writes
the merge and the test command passes.}

\begin{tikzpicture}[x=1mm, y=1mm, line cap=round, font=\sffamily]

% one function body. #1 x, #2 y of top edge, #3 mode: t / b / a
% t = the three lines at the top are gone.  b = a line added at the bottom.
\newcommand{\fbody}[3]{%
  \draw[card, color=sgtGrey!45, fill=white] (#1,{#2-9.4}) rectangle ({#1+13},#2);
  \foreach \yy in {-1.4,-2.6,-3.8,-5.0,-6.2,-7.4,-8.6}
    {\fill[sgtGrey!25] ({#1+1.3},{#2+\yy-0.3}) rectangle ({#1+9.4},{#2+\yy+0.15});}
  \if#3b\else
    \fill[white] ({#1+0.9},{#2-4.4}) rectangle ({#1+12.1},{#2-0.9});
    \draw[fCache, line width=0.3pt, densely dotted]
      ({#1+1.3},{#2-4.3}) rectangle ({#1+11.2},{#2-1.0});
  \fi
  \if#3t\else
    \fill[fIdx] ({#1+1.3},{#2-8.9}) rectangle ({#1+11.2},{#2-8.45});
  \fi}
% #1 x, #2 y of body top, #3 colour, #4 letter
\newcommand{\fbodytag}[4]{%
  \node[tag, fill=#3, text=white, anchor=south west,
        font=\sffamily\bfseries\fontsize{5.0}{5.8}\selectfont] at (#1,{#2+0.5}) {#4};}

% ============== 1: where the two versions come from ==================
\panel{0}{38}{84}{24}{sgtInk}{1\ \ ONE SAVED VERSION, TWO NEW ONES}

\draw[rail=sgtGrey] (2.5,48.5) -- (9.0,48.5);
\opext{5.5}{48.5}{sgtGrey}
\draw[raildim] (9.0,48.5) -- (13.5,53.0);
\draw[raildim] (9.0,48.5) -- (13.5,44.0);
\draw[rail=fCache] (13.5,53.0) -- (20.5,53.0); \oprew{17.0}{53.0}{fCache}
\draw[rail=fIdx]   (13.5,44.0) -- (20.5,44.0); \opext{17.0}{44.0}{fIdx}
\whodev{23.2}{53.0}{fCache}
\whodev{23.2}{44.0}{fIdx}

\node[anchor=west, text width=21mm, align=left, inner sep=0, text=fCache,
      font=\sffamily\fontsize{5.4}{6.2}\selectfont] at (26.4,53.2)
  {\textbf{A} the cache revert cut 3 lines off the top};
\node[anchor=west, text width=21mm, align=left, inner sep=0, text=fIdx,
      font=\sffamily\fontsize{5.4}{6.2}\selectfont] at (26.4,44.2)
  {\textbf{B} a colleague logged a request id at the end};

\fbody{53.0}{57.6}{t}\fbodytag{53.0}{57.6}{fCache}{A}
\fbody{68.5}{57.6}{b}\fbodytag{68.5}{57.6}{fIdx}{B}
\node[acc=sgtInk!70, anchor=north] at (67.2,46.8) {forty lines apart};

% ===================== 2: what git produces ==========================
\panel{0}{18}{84}{18}{sgtGrey}{2\ \ WHAT GIT PRODUCES}

\node[anchor=north west, text width=46mm, align=left, inner sep=0, text=sgtInk!85,
      font=\sffamily\fontsize{5.7}{6.7}\selectfont] at (2.2,29.8)
  {The two changes touch no common line, so git merges them and prints nothing.
   Nobody has run this function without its cache lookup and with the new log
   line.};

\fbody{53.0}{30.2}{a}
\fbodytag{53.0}{30.2}{sgtGrey}{A+B}
\annot{69.5}{24.4}{66.6}{25.4}{18}{fExport}
\node[anchor=west, text width=13mm, align=left, inner sep=0, text=fExport,
      font=\sffamily\bfseries\fontsize{5.4}{6.2}\selectfont] at (70.0,24.4)
  {a state nobody ran};

% ===================== 3: what sgt produces ==========================
\panel{0}{0}{84}{16}{sgtInk}{3\ \ WHAT \textsf{sgt} PRODUCES}

\node[anchor=north west, inner sep=0, text=sgtInk,
      font=\ttfamily\fontsize{5.6}{6.4}\selectfont] at (2.2,11.6)
  {\$ sgt resolve api.py::handle};

\draw[raildim] (2.4,5.6) -- (6.0,7.8);
\draw[raildim] (2.4,5.6) -- (6.0,3.4);
\draw[raildim] (6.0,7.8) -- (11.6,7.8); \opghost{8.8}{7.8}{fCache}
\draw[raildim] (6.0,3.4) -- (11.6,3.4); \opghost{8.8}{3.4}{fIdx}
\node[anchor=west, text width=20mm, align=left, inner sep=0, text=sgtInk!80,
      font=\sffamily\fontsize{5.4}{6.2}\selectfont] at (13.4,5.6)
  {both recorded,\\neither included};

\draw[-{Straight Barb[length=0.9mm,width=1.0mm]}, line width=0.35pt, color=sgtGrey]
  (35.0,5.6) -- (38.6,5.6);
\draw[rail=sgtInk] (40.6,5.6) -- (47.6,5.6); \oprew{44.1}{5.6}{sgtInk!70}
\whodev{50.4}{5.6}{sgtInk!70}
\node[anchor=west, text width=29mm, align=left, inner sep=0, text=sgtInk!85,
      font=\sffamily\fontsize{5.4}{6.2}\selectfont] at (53.0,5.6)
  {the version the developer wrote, included once the tests pass};

\end{tikzpicture}

  \Description{\FigForkDescription}
  \caption{\FigForkCaption}
  \label{fig:fork}
\end{figure}

When two edits change the same function from the same starting version, \sgt{}
records both, includes neither, and asks the developer. This is the commitment that
costs the most goodwill in the first week, so we want to be precise about whose
behaviour we are departing from and why theirs is right.

Git merges text and knows nothing about what a function is. A system that has
to work on every text file in every language anyone will ever write can afford to
know nothing else, which is why \cmd{git merge} is useful on a Makefile. In the case
Figure~\ref{fig:fork} draws, two people had last saved the same version of
\sym{api.py::handle}, one of them reverted a caching feature, which took the cache
lookup out of the first three lines, and the other added a request identifier to the
logging block forty lines away, with no line in common. Git merges the two without
comment, and the function that comes out has never been run in that state by
anybody. \sgt{} sees one function with two versions from one parent, names it, and
stops. \cmd{sgt resolve api.py::handle} prints the two versions beside each other
and asks \cmd{pick [l]eft / [r]ight / [e]dit / [q]uit}, where the third choice opens
the file in the developer's editor and the fourth leaves the fork open, and whatever
they write has to pass the project's test command before \sgt{} will include it.
Where no test command is configured the developer supplies the verdict with
\cmd{--override pass} and a \cmd{--reason}, and the reason and their name are
recorded beside the resolution, so the next person to read that function finds out
why it looks the way it does.

A developer pays for this by being stopped in cases where the merge would have been
fine, and in one of those cases \sgt{} stops for nothing at all, since deleting a
function and adding it back with byte-identical content in two places produces two
versions from one parent and \sgt{} treats content that agrees exactly as a
disagreement. The design is a fail-stop commitment in Schlichting and Schneider's
sense~\cite{schlichting1983}: the system converts a potentially-wrong state (a
textual merge whose correctness nobody has verified) into a detectable halt (the
developer is asked). Fail-stop costs availability---the tool is unavailable until the
developer answers---in exchange for a guarantee about everything it does let through:
if the system did not stop, the included version is one a person chose, not one a
program composed. The read layer inherits this property: every attribution,
dependency display, and consequence preview rests on versions that entered through a
human decision, so the reports are as trustworthy as the decisions they trace back to.
A silent textual merge would break that chain, because the function it produces has
never been chosen by anybody, and the next revert that touches it would be operating on
text whose provenance is a program's guess rather than a person's act.

We hold the commitment loosely, because a
developer who merges thirty branches a day would reasonably weigh the two the other
way, and for them the current design is worse. The 18$\times$ amplification
Section~\ref{sec:session-rate} reports is the measured cost of choosing correctness
over availability at this grain, and a tool builder facing the same choice should
treat that number as a lower bound on what fail-stop costs at the function level.

The deeper reason for accepting that cost is the context in which the tool
operates. A developer who typed the competing versions themselves would have read
both, would hold a model of what each does, and could evaluate a textual merge in
seconds. A developer whose agent produced one version while a colleague's agent
produced the other has read neither, holds no model of the interaction between them,
and would accept a textual merge not because they verified it but because the tool
offered no reason to look. Shen and Tamkin measured this directly: developers who
relied on AI assistance scored 17\% lower on assessments of the code they had just
produced, with debugging---recognising when code is wrong---showing the largest
gap~\cite{shen2026skills}. The fork rule is designed for the developer in that
degraded state: it forces a decision point that requires looking at the two versions,
which is the cognitive engagement their earlier delegation removed. A tool that
silently merges code its user has not read, in a context where the user's capacity to
read critically has already declined, is compounding the problem it was built to
address.

\subsection{The names are inferred, and correcting one moves labels only}
\label{sec:names}

Every figure in this paper labels a group of edits with a name like
\emph{price cache}, and those names are the part of \sgt{} that guesses. Functions
that changed in the same saves and that refer to one another are grouped together,
using the change coupling signal established by Zimmermann et al.\ and Ying et
al.~\cite{zimmermann2004,ying2004} and a community detection method over the
resulting graph~\cite{traag2019}. The reason a graph algorithm finds real features
rather than arbitrary cuts is that codebases exhibit what Simon called
\emph{near-decomposability}: interactions within a subsystem are strong and frequent
while interactions between subsystems are weak and
infrequent~\cite{simon1962}. Community detection recovers that structure because
it was designed to, and the known exception---a hub function called from
everywhere---is the same entity that breaks feature location in the literature and
that Section~\ref{sec:discussion-synthesis} reports as the main source of
grouping errors here: it pulls otherwise separate requests into one community
because the coupling signal cannot tell structural proximity from shared purpose.

A language model then reads the saves in a group
and writes the label, which is the descendant of a task attempted before language
models existed, when Buse and Weimer generated a description of a change out of the
change itself~\cite{buse2010}. The label draws from the developer's words rather than
from the code's identifier names. Nguyen and
Glassman showed that the naming a model uses in its output and the naming a
developer uses to describe the same behaviour systematically
mismatch~\cite{nguyen2024}. A label derived from identifier names would therefore speak the
model's language rather than the developer's, and the developer who later searches
for ``the caching'' would find nothing under whatever the model called its cache
class. Membership is computed from the code and the label is written from the words,
and those two steps fail in different ways, so we treat them separately. The
grouping is also systematically finer than the request: a community detection
algorithm cuts where the coupling is weak, not where the developer stopped
typing, so one request that touches several disjoint parts of the code becomes
several features. Section~\ref{sec:discussion-synthesis} reports a median of
4--6 features per request across four repositories. The result is a sub-intent
decomposition rather than an intent recovery---the developer gets a vocabulary
for the structural neighbourhoods their request produced, not a map back to what
they asked for. The finer grain has a cost and a benefit. The cost is that the
developer cannot name what they asked for in one selection; they must merge or
select several. The benefit is that operations work at a finer level than the
request: the developer who asked for caching and later wants only the TTL
helper removed can name that piece without reverting the whole request. The
mismatch is in level as well as in number: the
developer thinks at the level of ``the price cache'' (one purpose, one name) while
the tool presents four structural communities inside it, a difference
Section~\ref{sec:discussion-synthesis} traces to Rosch's distinction between
basic-level and subordinate categories~\cite{rosch1978}.

We chose to operate at the subordinate level rather than the basic level for a
reason that becomes clear only when the developer acts on the grouping. A
request boundary is unstable: it depends on what the developer typed, which changes
between ``implement caching'' (one request) and ``add a TTL helper, then wire it
into the handler'' (two requests for the same code). A structural boundary is
stable: it depends on which functions call which others, which does not change when
the developer rephrases. A revert, a restore, or a dependency warning all need a
boundary that means the same thing regardless of when the developer asks---and a
boundary defined by coupling has that property while a boundary defined by request
phrasing does not. The over-decomposition is therefore not an error in the
clustering but a consequence of choosing stability over correspondence: the tool
gives four addressable pieces where the developer expected one, because those four
pieces are the ones whose membership will not shift when the next request arrives.
The cost (merging to act at the request level) is paid once; the benefit (stable
addresses for operations) pays on every subsequent use.

We considered asking the developer to confirm the grouping at save time and decided
against it, because answering that question takes the same knowledge commit hygiene
takes, and a developer who could reliably answer it could also have staged the
hunks. Shipman and Marshall showed that systems which ask users to formalise their
own work get abandoned~\cite{shipman1999}; the design they arrived at has the system
propose a structure the user may accept, modify, reject, or leave alone, which is
the arrangement here. Our conditions are easier than theirs: the structure their
systems tried to elicit was tacit, so a user had first to work out what they thought.
The structure \sgt{} recovers was spoken in a request before any of the code
existed---though not at the grain the tool produces, since one request typically
splits into several structural clusters. The reason it needs recovering at all
is that nothing wrote it down. What matters is therefore not whether the guess is
right but what a wrong one costs, and correcting one moves labels and touches no
recorded edit. \cmd{sgt feature rename 001f73c2 "price cache"} renames a group whose
identifier stays what it was, \cmd{sgt feature regroup} moves functions between
groups, and \cmd{sgt save --as "price cache"} files a new save into a group that
already exists. Because identifiers survive all three, a revert written down in a
code review comment last month still removes the same edits after somebody renames
the group today.

The deeper function of inferred names is coordination. Clark's analysis of language
use shows that people who need to refer to the same thing build a shared
vocabulary through proposal and acceptance: one party offers a term, the other
either adopts it or repairs it, and after a few exchanges the term is common
ground~\cite{clark1996}. A version control system that infers feature names is
proposing one half of that exchange---offering ``price cache'' as the term both the
developer and the tool will use to pick out a set of edits. If the developer types
\cmd{sgt revert "price cache"} the proposal has been accepted and the name is
grounded; if they rename it, the repair has completed the same grounding through a
different path. Either way, the tool and developer now share a vocabulary for the
codebase, and that vocabulary enters the developer's prompts to the AI assistant,
which is why the ``prompt specificity'' measure in Section~\ref{sec:findings-process}
tests whether the tool's names reached the developer's language. A wrong label costs
a rename (one command, no data moved); a missing shared vocabulary costs every
subsequent conversation about the code, because the developer and the tool have no
common referent for the set of edits they both need to name.

Two failures show up often enough that a reader will hit them. A function that
everything calls, an argument parser or a logging helper, changes in nearly every
save and pulls otherwise unrelated work into one group, and the fix is a regroup the
developer has to notice they need. A save with no message gives the language model
less to read, and \sgt{} says so at the time, printing \cmd{\textperiodcentered{} no words captured}
under the save it just wrote, because the label summarises what the developer said
about the work and a developer who says nothing gets a name derived from the code
alone.

\subsection{What \sgt{} will not do}

\sgt{} writes no code of its own. Every byte in the working files came from the
developer or their agent, apart from the three-way merge splice in
Section~\ref{sec:sets}, which reports every case it cannot perform cleanly. We have
been asked more than once during this work why \sgt{} does not hand the shared
function to a language model and have it redrafted after a revert. We considered it
seriously and rejected it for one reason: a revert whose output depends on a
model's judgement produces a file whose correctness the developer can verify only
by reading every line the model touched, which is the labour Section~\ref{sec:scenarios}
argued this tool should eliminate. A revert whose output is computed from the
records the developer already validated produces a file whose correctness follows
from the records, and the developer's test suite decides whether the composition
is sound. The cost is that a subtraction which overlaps later work requires a hand
edit; the benefit is that the developer knows what they are reading. Bansal et
al.\ found that AI explanations increase acceptance whether or not the underlying
recommendation is correct~\cite{bansal2021}; a model-in-the-loop revert would
face the same problem, because the developer would have to judge whether the
model's redrafted function is correct, and the tool's report that it passed tests
is exactly the kind of explanation that increases acceptance without increasing
accuracy. The refusal has a deeper justification. A version control system is
infrastructure: it sits beneath every other decision, and a revert that is wrong
invalidates every test, every review, and every deployment that followed it. An
AI assistant is a collaborator: its suggestions are proposals that a developer
accepts, modifies, or rejects one at a time, and a wrong suggestion costs one
undo. The failure modes are not symmetric. A wrong line completion costs the
developer a delete; a wrong revert costs them their codebase state, and the
difference in blast radius is the difference between a tool that should reach for
a model (the assistant) and one that should not (the substrate). The trend toward
model-in-the-loop for every operation conflates the two roles, and the
conflation is invisible as long as the model is correct. The moment it is wrong,
the developer discovers that the substrate contained a guess they never verified,
and the cost of discovering it late is what Section~\ref{sec:discussion}
measures.

\sgt{} also does not stand between a developer and git. The records are files under
\cmd{.sgt/} committed alongside the code, \cmd{sgt git <args>} passes through to git
untouched, a colleague who has never installed \sgt{} clones an ordinary repository
with an ordinary history, and there is no server and no background process. Two
commands refuse to run: \cmd{switch} stops when there are unsaved edits, because
\sgt{} has nowhere to put them, and names the files and prints \cmd{record them with
`sgt save', or commit / `git restore' those files, then re-run}; \cmd{restore} stops
when putting a set back would leave two versions of one function live, and there its
message names the edits by identifier while the function whose two versions caused
the refusal goes unnamed, which is the rough edge Section~\ref{sec:walkthrough}
reports us fixing in the fork warning and not here. Everything else
\sgt{} finds wrong it reports, and then it proceeds.

\subsection{From the five operations to the commands a developer types}

\begin{figure*}[t]
  \centering
  % Every operation sgt offers, grouped by what it does to the recorded set, with
% the state before and after each command and what each one costs.
% Prints as Figure 5, after the fork figure in Section 4.3.
\newcommand{\FigFourCaption}{The nine operations that act on the record, grouped
by what each one does to it. One command writes records and the other eight read
them or change which of them a version includes, and no argument anywhere in the
figure is a hash or a line number: every one of them is a request the developer
named, a branch, or a function.}

\newcommand{\FigFourDescription}{Three panels, one per group of operations. The
first panel holds the single operation that writes records, the second holds the
three that only read them, and the third holds the five that change which
records a version includes. Each row carries a small diagram of the set before
and after the command, the command itself, one sentence saying what the developer
gets, and one sentence in a grey band on the right saying what the operation
costs. Each diagram draws one horizontal line per request with a mark for every
edit filed under that request, as in Figure 1, and a dotted outline is an edit
the version no longer includes.}

\begin{tikzpicture}[x=1mm, y=1mm, line cap=round, font=\sffamily]
% narrow ragged columns, so break a line early rather than hyphenate a word
\hyphenpenalty=10000 \exhyphenpenalty=10000

% the thumbnails carry the before/after convention, so the column says so itself
\node[anchor=west, inner sep=0, text=sgtInk!60,
      font=\sffamily\fontsize{5.0}{5.8}\selectfont] at (3.4,104.6)
  {before \ttarrow{} after};
\node[hd, anchor=west] at (19,104.6)    {COMMAND};
\node[hd, anchor=west] at (53,104.6)    {WHAT THE DEVELOPER GETS};
\node[hd, anchor=west] at (115.5,104.6) {WHAT IT COSTS};

% one group of operations. The cost column gets a band of its own on the right,
% the same band Figure 2 uses, so the reader can read that column straight down.
% #1 y, #2 height, #3 tab colour, #4 tab text
\newcommand{\opgroup}[4]{%
  \panel{0}{#1}{176}{#2}{#3}{#4}
  \fill[wGrey, rounded corners=0.8mm] (112.5,{#1+1.2}) rectangle (174.6,{#1+#2-1.2});}

% #1 y of the first text line, #2 thumbnail, #3 command, #4 gets, #5 costs
\newcommand{\oprow}[5]{%
  \node[inner sep=0] at (10,{#1-2.1}) {#2};
  \node[anchor=north west, text width=32mm, align=left, inner sep=0, text=sgtInk,
        font=\ttfamily\fontsize{5.6}{6.6}\selectfont] at (19,#1) {#3};
  \node[anchor=north west, text width=55mm, align=left, inner sep=0, text=sgtInk!85,
        font=\sffamily\fontsize{5.7}{6.7}\selectfont] at (53,#1) {#4};
  \node[anchor=north west, text width=56mm, align=left, inner sep=0, text=sgtInk!85,
        font=\sffamily\fontsize{5.7}{6.7}\selectfont] at (115.5,#1) {#5};}

% =================== what writes records ==================================
% One row, and the group holding one row is the point: the developer types one
% command to add to the record, and everything else in the figure decides what
% that record is read for or what a version does with it.
\opgroup{86.4}{13.6}{sgtInk}{WHAT WRITES RECORDS}

% Centred in its panel rather than hung from the top of it, so that the space
% around the one row reads as emphasis and not as a row somebody took out.
\oprow{94.3}{\thsave}{sgt save -m "..."}
  {One record per changed function, filed under a name for the request. The only
   command that adds to the record.}
  {The grouping is a guess. Correcting it moves labels and changes no recorded
   edit.}

% =================== what only reads them =================================
\opgroup{55.4}{29}{sgtGrey}{WHAT ONLY READS THEM}

\oprow{79.8}{\thwhy}{sgt show api.py::handle}
  {The request this function came from, the saves that touched it, and how much
   other work removing it takes along.}
  {The answer is exactly as good as the grouping behind it.}
\oprow{71.4}{\thdiff}{sgt diff main feat/export}
  {The edits each version has that the other does not, in both directions.}
  {It answers in functions, so seeing which lines moved still means opening the
   file.}
\oprow{63.0}{\thckpt}{sgt intent list}
  {Each feature cut into the stretches of work done on it, each with a name like
   {\ttfamily rate limiting@2} the other commands accept.}
  {A language model places the cuts by reading the saves, so a stretch can start
   one save too early.}

% ============ what changes which records a version includes ===============
\opgroup{7.4}{46}{sgtInk}{WHAT CHANGES WHICH RECORDS A VERSION INCLUDES}

\oprow{48.8}{\threvert}{sgt revert "price cache"}
  {The named set subtracted from the version, with every function it changes
   named before anything is applied.}
  {A function later work also edited gets the set's lines subtracted from its
   current text, and the tests decide.}
\oprow{40.4}{\threstore}{sgt restore "price cache"}
  {The same set back, byte for byte, with everything added since left alone.}
  {Refused when putting the set back would leave two versions of one function
   live, and the refusal prints edit identifiers with no function name.}
\oprow{32.0}{\thswitch}{sgt switch feat/export}
  {The working files rebuilt from that branch's set.}
  {Refused on unsaved edits, because \sgt{} has nowhere to put them yet.}
\oprow{23.6}{\thfork}{sgt resolve api.py::handle}
  {The two competing versions side by side, and whichever one the developer
   picks or writes is included once the project's test command passes.}
  {Two edits at opposite ends of one function are also competing versions, and
   git would have merged them without asking.}
\oprow{15.2}{\thsubset}{sgt propose land --subset "export"}
  {One finished feature sent to the shared branch, with the unfinished one left
   behind.}
  {Two functions that sit next to each other in one file can lose the blank line
   between them.}

\end{tikzpicture}

  \Description{\FigFourDescription}
  \caption{\FigFourCaption}
  \label{fig:operations}
\end{figure*}

Figure~\ref{fig:operations} sets out the nine operations that act on the record,
each with its cost beside it, and the five operations Section~\ref{sec:scenarios}
ended on are four commands and one thing that needs no command. Recording a
boundary without declaring it is
\cmd{sgt save -m}, run whenever the developer feels like running it, since the
grouping is recovered afterwards from the code. Seeing the work one request produced
is \cmd{sgt show "rate limiting"}, or \cmd{sgt show "rate limiting@1"} for one
stretch of work on it alone. Removing a named set is
\cmd{sgt revert "price cache"}, and carrying one to another branch is
\cmd{sgt propose land --subset "export"}, which declines a subset that omits a group
another chosen group depends on and names the group that is missing. Each of those
four takes as its argument a phrase somebody used while the work was happening, so
the manual recovery Section~\ref{sec:scenarios} described is replaced by a named
command---on repositories where the store reproduces the committed bytes. Where
it does not (27 of 28 outside repositories), the command refuses rather than
producing a wrong file, and the developer is back to the manual recovery with no
help from the tool (Section~\ref{sec:session-rate}).

Discovering the
layering needs no command at all wherever the references resolved, because the cache
edits record that they were built on \sym{api.py::\_bucket} and therefore on the rate
limiter, so the developer has at any later moment the fact they needed at 14:30 and
did not have. Where they did not resolve, the layering is simply absent rather than wrong, and the
developer is back to reading the code, which is the position the ambiguity rule
above trades down to on purpose. On 35 corpus repositories, 93\% of ops carry no
reference at all, because resolution requires a globally unique function name---so
the absence is the common case, not a rare failure. The 7\% that do resolve tend to
be cross-module calls (a handler calling a cache function in another file, a test
importing a fixture), which are precisely the dependencies whose absence would break
the code silently---so the coverage is thin but concentrated on the cases where a
human would not have noticed the dependency either. Where resolution is absent, the
developer's test suite is the only oracle, and \sgt{} says so in the revert output
(``still references removed code'') only when it can, which on most repositories is
never. The implication for adoption is blunt: a developer whose repository has few
cross-module calls will never see a dependency warning, and will discover that the
tool offers no protection there only when a revert breaks something silently.

A reader should treat the dependency display as a bonus that fires
occasionally rather than a guarantee the design rests on; the design rests on set
arithmetic over edits, and the dependency enriches a revert without being required
for one. The rest of
what a developer types is housekeeping around those nine, \cmd{sgt init} to begin,
\cmd{sgt log} and \cmd{sgt sync} to see where things stand, and the
\cmd{sgt feature rename} of Section~\ref{sec:names} to correct a name, none of which
changes which edits a version includes. Section~\ref{sec:walkthrough} runs these in
order on one repository.
